\documentclass{beamer}
\usetheme{Berlin}

\usepackage{xcolor}
\usepackage{listings}

\lstset{
    language=C++,
    basicstyle=\ttfamily\small,
    keywordstyle=\color{blue},
    commentstyle=\color{green!50!black},
    stringstyle=\color{red},
    showstringspaces=false,
    frame=single,
    breaklines=true,
    columns=fullflexible
}

\title{Compile-Time Sparse Vector Expressions}
\subtitle{Expression Templates with Projection}
\author{Christian Kotz}
\date{}

\begin{document}

\frame{\titlepage}

%--------------------------------------------------

\begin{frame}{Motivation}

\begin{itemize}
\item Sparse data structures appear everywhere
\item Classical evaluation computes \emph{everything}
\item Often we only need a \emph{small subset} of indices
\end{itemize}

\vspace{1em}

\textbf{Goal:}
\begin{itemize}
\item Build expressions lazily
\item Evaluate only the indices we actually need
\end{itemize}

\end{frame}

%--------------------------------------------------

\begin{frame}[fragile]{Expression Templates (recap)}

\begin{lstlisting}
auto expr = a + b - c;
\end{lstlisting}

\vspace{1em}

Builds an AST:
\begin{itemize}
\item nodes: operations (+, -)
\item leaves: values
\end{itemize}

\vspace{1em}

\textbf{Idea:} delay evaluation

\end{frame}

%--------------------------------------------------

\begin{frame}[fragile]{Key Idea: Projection}

\begin{lstlisting}
using I = ctv::index_set<1,4,5,9>;

auto result = evaluate<I>(a + b - c);
\end{lstlisting}

\vspace{1em}

\textbf{Meaning:}

\begin{itemize}
\item Only indices in $I$ are evaluated
\item Everything else is treated as zero
\end{itemize}

\vspace{1em}

\textbf{Interpretation:}
\begin{itemize}
\item evaluation = projection onto a support set
\end{itemize}

\end{frame}

%--------------------------------------------------

\begin{frame}[fragile]{Evaluation = Projection on AST}

\begin{columns}[T]
\begin{column}{0.45\textwidth}
\textbf{AST}

\begin{verbatim}
      (-)
     /   \
   (+)    c
  /   \
 a     b
\end{verbatim}
\end{column}

\begin{column}{0.5\textwidth}
\textbf{Projection set}

\[
I = \{1,4,5,9\}
\]

\vspace{0.5em}

\textbf{Queried indices}

\[
\begin{aligned}
1 &\mapsto eval\_at(1, expr) \\
4 &\mapsto eval\_at(4, expr) \\
5 &\mapsto eval\_at(5, expr) \\
9 &\mapsto eval\_at(9, expr)
\end{aligned}
\]
\end{column}
\end{columns}

\vspace{0.8em}

\textbf{Key idea:} We do not traverse the tree --- we query it.

\end{frame}

%--------------------------------------------------

\begin{frame}[fragile]{Terminal: sparse\_vector}

\begin{lstlisting}
template <typename IndexSet, typename T>
struct sparse_vector
{
    using index_set = IndexSet;
    using value_type = T;

    std::array<T, IndexSet::size> values;
};
\end{lstlisting}

\vspace{1em}

\begin{itemize}
\item compile-time support set
\item dense storage over sparse indices
\item also acts as expression terminal
\end{itemize}

\end{frame}

%--------------------------------------------------

\begin{frame}[fragile]{Expression Nodes}

\begin{lstlisting}
template <typename LHS, typename RHS>
struct add_expr
{
    LHS lhs;
    RHS rhs;
};
\end{lstlisting}

\vspace{1em}

\begin{lstlisting}
auto expr = a + b;
\end{lstlisting}

\vspace{1em}

\begin{itemize}
\item builds AST
\item no computation yet
\end{itemize}

\end{frame}

%--------------------------------------------------

\begin{frame}[fragile]{Evaluation per Index}

\begin{lstlisting}
eval_at(i, expr)
\end{lstlisting}

\vspace{1em}

Recursive:

\begin{lstlisting}
eval_at(i, a + b)
= eval_at(i, a) + eval_at(i, b)
\end{lstlisting}

\vspace{1em}

Terminal:

\begin{lstlisting}
eval_at(i, sparse_vector)
\end{lstlisting}

returns value or 0

\end{frame}

%--------------------------------------------------

\begin{frame}[fragile]{Materialization}

\begin{lstlisting}
evaluate<I>(expr)
\end{lstlisting}

\vspace{1em}

\begin{itemize}
\item iterate over compile-time index set
\item compute each entry via code \lstinline!eval_at!
\item store in array
\end{itemize}

\vspace{1em}

\textbf{Key point:}
\begin{itemize}
\item work is proportional to requested indices
\end{itemize}

\end{frame}

%--------------------------------------------------

\begin{frame}[fragile]{Full Support Set}

\begin{lstlisting}
full_index_set_t<Expr>
\end{lstlisting}

\vspace{1em}

\begin{itemize}
\item computed from AST
\item union of all terminal supports
\end{itemize}

\vspace{1em}

\begin{lstlisting}
evaluate(expr)
\end{lstlisting}

\vspace{1em}

\begin{itemize}
\item evaluates on natural support
\end{itemize}

\end{frame}

%--------------------------------------------------

\begin{frame}[fragile]{Example}

\begin{lstlisting}
a: {1,3,5}
b: {3,4,5}

expr = a + b
\end{lstlisting}

\vspace{1em}

Full support:

\begin{itemize}
\item {1,3,4,5}
\end{itemize}

\vspace{1em}

\begin{lstlisting}
evaluate<{1,5}>(expr)
\end{lstlisting}

\begin{itemize}
\item only computes 1 and 5
\end{itemize}

\end{frame}

%--------------------------------------------------

\begin{frame}{Design Decision}

\textbf{Operands stored by value}

\vspace{1em}

\begin{itemize}
\item safe for temporaries
\item avoids lifetime bugs
\item simpler mental model
\end{itemize}

\vspace{1em}

Production approach:

\begin{itemize}
\item lvalues $\rightarrow$ reference
\item rvalues $\rightarrow$ move
\end{itemize}

\end{frame}

%--------------------------------------------------

\begin{frame}{Summary}

\begin{itemize}
\item Expression templates build an AST
\item Evaluation is delayed
\item We add: evaluation restricted by index set
\end{itemize}

\vspace{1em}

\textbf{Core idea:}

\begin{itemize}
\item evaluation = projection onto support set
\end{itemize}

\vspace{1em}

\textbf{Result:}

\begin{itemize}
\item compute only what is needed
\end{itemize}

\end{frame}

%--------------------------------------------------

\begin{frame}{Outlook}

\begin{itemize}
\item context-driven evaluation
\item symbolic variables
\item Clifford / geometric algebra
\item index-dependent operations
\end{itemize}

\vspace{1em}

\textbf{Key idea remains:}

\begin{itemize}
\item projection + context
\end{itemize}

\end{frame}

\end{document}
